1. Ce trebuie să calculăm?
Alex scoate, fără înlocuire, K șosete dintr-un sertar în care fiecare șosetă are

o culoare (numere de la 1 la N) și
o dimensiune – mică (size = 0) sau mare (size = 1).
Toate șosetele mici au aceeași greutate sizeSmall, toate cele mari au greutatea sizeBig.
În fiecare pas, șansele ca o anumită șosetă să fie aleasă sunt proporționale cu greutatea ei.

După ce a luat cele K șosete, Alex vrea să ştie probabilitatea ca cel puţin două dintre ele să aibă aceeaşi culoare (dimensiunea nu contează).
Se cere răspunsul cu eroare absolută ≤ 10⁻⁵.

2. Ideea principală
Evenimentul „există cel puţin o pereche de şosete de aceeaşi culoare” este complementul evenimentului „toate şosetele alese au culori diferite”.

[ P_{\text{cel puţin două egale}} = 1 - P_{\text{toate diferite}} . ]

Așadar trebuie să calculăm probabilitatea ca, în cele K alegeri, să nu se repete nicio culoare.

Greutatea depinde numai de dimensiune, nu de culoare. Din acest motiv procesul poate fi descompus în două părţi independente:

Câte şosete mici și câte mari sunt luate?
Dacă notăm cu a numărul de mici și cu b = K‑a numărul de mari, distribuţia perechii (a,b) se obţine printr-un DP simplu (q[a][b]).

Pentru o pereche fixă (a,b), care este probabilitatea ca toate culorile să fie distincte?
Odată cunoscut că exact a mici și b mari au fost scoase, selecţia dintre şosetele mici (respectiv mari) este uniformă: fiecare submulţime de a elemente ale celor A mici și fiecare submulţime de b elemente ale celor B mari este la fel de probabilă. Astfel problema devine pur combinatorială – numărarea submulţimilor care conţin cel mult un şosetă din fiecare culoare.

Răspunsul final se obţine prin sumarea peste toate combinaţiile posibile (a,b):

[ \boxed{\displaystyle \text{Răspuns}= \sum_{a=0}^{K} q[a][K-a];\bigl(1-p[a][K-a]\bigr) } ]

unde

q[a][b] – probabilitatea ca în cele K alegeri să existe exact a şosete mici și b şosete mari;
p[a][b] – probabilitatea, condiţionată de existenţa acelei repartizări de dimensiuni, că toate culorile sunt diferite.
3. Calculul distribuţiei q[a][b]
Să notăm

A – numărul total de şosete mici,
B – numărul total de şosete mari,
și greutăţile lor wS = sizeSmall, wB = sizeBig.

Definim

[ q[i][j] = \Pr\bigl(\text{după }i+j\text{ alegeri am luat }i\text{ mici şi }j\text{ mari}\bigr). ]

Starea iniţială este q[0][0] = 1.
Dintr‑o stare (i,j) greutatea totală rămasă este

[ W = (A-i) wS + (B-j) wB . ]

Probabilitatea de a alege o şosetă mică (dacă mai există) este

[ \frac{(A-i) wS}{W}, ]

iar cea de a alege o şosetă mare este

[ \frac{(B-j) wB}{W}. ]

Prin urmare tranziţiile DP sunt

[ \begin{aligned} q[i+1][j] &\mathrel{+=} q[i][j];\frac{(A-i) wS}{(A-i) wS+(B-j) wB},\ q[i][j+1] &\mathrel{+=} q[i][j];\frac{(B-j) wB}{(A-i) wS+(B-j) wB}. \end{aligned} ]

Se completează tabelul pentru i = 0 … A și j = 0 … B.
Complexitatea este O(A·B) ≤ O(N²) – neglijabilă pentru N ≤ 199.

4. Probabilitatea ca toate culorile să fie diferite pentru un (a,b) fix
4.1 Datele pe culori
Pentru fiecare culoare c numărăm

cntS[c] – câte şosete mici de acea culoare există,
cntB[c] – câte şosete mari de acea culoare există.
Numărăm doar culorile care apar efectiv; să le notăm 1 … M (M ≤ N).

4.2 Ce putem lua dintr‑o culoare?
Pentru a nu repeta nicio culoare, dintr‑o culoare c putem lua cel mult una şosetă, iar aceasta poate fi:

niciuna – 1 posibilitate,
o mică – există cntS[c] moduri de a alege una dintre ele,
o mare – există cntB[c] moduri de a alege una dintre ele.
Nu putem lua simultan o mică și o mare din aceeași culoare, pentru că ar crea deja o repetiție.

4.3 DP care numără submulţimile valide
Definim

[ dp[i][x][y] = \text{numărul de moduri de a selecta }x\text{ mici şi }y\text{ mari din primele }i\text{ culori, fără repetiţii de culoare}. ]

Iniţial dp[0][0][0] = 1, restul 0.

Pentru fiecare culoare i = 1 … M se aplică cele trei alegeri posibile:

nu luăm nimic → dp[i][x][y] += dp[i-1][x][y];
luăm o mică (dacă există) → dp[i][x][y] += dp[i-1][x-1][y] * cntS[i];
luăm o mare (dacă există) → dp[i][x][y] += dp[i-1][x][y-1] * cntB[i].
După procesarea tuturor culorilor, dp[M][a][b] este exact numărul de submulţimi care conţin a mici și b mari și în care nicio culoare nu apare de două ori.

4.4 Normalizarea
Dacă ignorăm condiţia de culoare, numărul de moduri de a alege a mici și b mari este pur combinatorial:

[ \text{tot}[a][b] = \binom{A}{a},\binom{B}{b}. ]

Prin urmare probabilitatea ca, în condiţia de a avea exact a mici şi b mari, toate culorile să fie diferite este

[ p[a][b] = \frac{dp[M][a][b]}{\text{tot}[a][b]} . ]

5. Combinarea rezultatelor
Pentru fiecare a = 0 … K (și b = K‑a) cunoaştem deja

q[a][b] – probabilitatea distribuţiei de dimensiuni (secţiunea 3);
p[a][b] – probabilitatea de a nu avea culori repetate (secţiunea 4).
Răspunsul cerut este

[ \text{Răspuns}= \sum_{a=0}^{K} q[a][K-a];\bigl(1-p[a][K-a]\bigr). ]

6. Demonstrarea corectitudinii
Lema 1
Condiționat de evenimentul „am luat exact a şosete mici și b şosete mari”, mulțimea de şosete mici luate este uniform aleatoare din toate submulțimile de a elemente ale celor A şosete mici; același lucru este valabil pentru şosetele mari.

Demonstrație.
Toate şosetele mici au aceeași greutate sizeSmall. În fiecare pas raportul dintre probabilitățile de a alege oricare două şosete mici este 1. Astfel procesul de extragere fără înlocuire este echivalent cu alegerea uniformă a unui subset de dimensiune a. Argumentul se aplică în mod simetric și pentru şosetele mari. ∎

Lema 2
Valorile q[a][b] obținute prin DP-ul din secțiunea 3 coincid cu probabilitatea reală ca în cele K alegeri să existe exact a şosete mici și b şosete mari.

Demonstrație.
DP‑ul urmărește exact legea probabilității totale: dintr‑o stare (i,j) se trece la (i+1,j) cu probabilitatea proporțională cu greutatea totală a şosetelor mici rămase și la (i,j+1) cu probabilitatea proporțională cu greutatea totală a şosetelor mari rămase. Prin inducție pe numărul de pași, toate căile care duc la (a,b) sunt contorizate cu probabilitatea corectă, iar suma lor este q[a][b]. ∎

Lema 3
„dp[M][a][b] este numărul exact de submulțimi ce conțin a șosete mici și b șosete mari și în care nicio culoare nu apare de două ori.”

Demonstratie
Fie C₁ , C₂ , … , C_M lista tuturor culorilor care apar în sertar (după comprimare).
Pentru fiecare culoare C_i cunoaștem două valori:

cntS[i] – numărul de șosete mici de acea culoare,
cntB[i] – numărul de șosete mari de acea culoare.
În construcția submulțimii dorite trebuie să respectăm două restricții stricte:

Numărul total de mici este exact a și al de mari este exact b.
Din fiecare culoare se poate alege cel mult o singură șosetă (oricare dintre cele mici, dintre cele mari, sau niciuna).
Aceste restricții pot fi exprimate în mod recursiv, procesând culorile una câte una.

Pasul de inducție
Presupunem că pentru primele i‑1 culori am calculat corect toate posibilitățile de a obține un anumit număr de mici și mari, fără a repeta nicio culoare. Aceste valori sunt conținute în tabelul dp[i‑1][x][y], unde x și y sunt numerele de mici, respectiv mari, alese până în acel moment.

Nu luăm nicio șosetă din C_i.
În acest caz numărul de mici și mari rămâne (x , y). Toate variantele deja contorizate în dp[i‑1][x][y] rămân valide, deci contribuim cu dp[i‑1][x][y].

Luăm o șosetă mică din C_i.
Această alegere este posibilă numai dacă cntS[i] > 0 și dacă mai avem nevoie de încăcel puțin un mic (x ≥ 1). Pentru fiecare dintre cele cntS[i] șosete mici disponibile există o modalitate distinctă de a o alege, iar restul alegerilor (pentru primele i‑1 culori) trebuie să producă exact x‑1 mici și y mari. Prin urmare adăugăm
dp[i‑1][x‑1][y] * cntS[i].

Luăm o șosetă mare din C_i.
Analog, dacă cntB[i] > 0 și y ≥ 1, pentru fiecare dintre cele cntB[i] mari există o modalitate de a o alege, iar restul alegerilor trebuie să producă x mici și y‑1 mari. Contribuția este
dp[i‑1][x][y‑1] * cntB[i].

Aceste trei cazuri sunt exhaustive și mutu exclusive:

Exhaustive – dintr‑o culoare nu există alte posibilităţi (nu putem alege simultan o mică și o mare, deoarece ar încălca condiţia de „nici o culoare nu apare de două ori”).
Mutu exclusive – fiecare submulțime respectă exact una dintre cele trei descrieri, deci nu există dublări în numărare.
Prin urmare, pentru fiecare pereche (x , y) relaţia de recurență

dp[i][x][y] = dp[i‑1][x][y]                     // nu luăm nimic
            + dp[i‑1][x‑1][y] * cntS[i] (dacă x>0) // luăm o mică
            + dp[i‑1][x][y‑1] * cntB[i] (dacă y>0) // luăm o mare
construiește exact numărul de submulțimi valide care includ primele i culori.

Pornind de la baza dp[0][0][0] = 1 (singura submulțime a niciunei culori este submulțimea vidă) și aplicând recurenţa pentru i = 1 … M, obținem prin inducție că, la final, dp[M][a][b] este exact numărul de submulțimi care conțin exact a șosete mici și b șosete mari și în care fiecare culoare apare cel mult o singură dată. ∎

7. Calculul probabilităţii finale
Pentru fiecare repartizare posibilă (a , b) cu a + b = K avem:

q[a][b] – probabilitatea ca în cele K extrageri să se obţină exact a mici și b mari (obţinută în prima parte a algoritmului);
p[a][b] = dp[M][a][b] / ( C(A,a)·C(B,b) ) – probabilitatea, condiționată de acea repartizare, ca toate culorile să fie diferite (normalizarea cu numărul total de combinaţii fără restricție de culoare).
Evenimentul „cel puţin două şosete au aceeaşi culoare” este complementul evenimentului „toate culorile sunt diferite”. Prin legea totală a probabilităţilor:

[ \begin{aligned} P(\text{cel puţin două egale}) &= 1 - \sum_{a+b=K} q[a][b];p[a][b] \ &= \sum_{a+b=K} q[a][b];\bigl(1-p[a][b]\bigr). \end{aligned} ]

Această formulă este exact cea utilizată în program. Suma are cel mult K+1 termeni, deci este calculată în timp trivial.

8. Complexitatea și utilizarea memoriei
DP pentru distribuţia de dimensiuni – O(A·B) timp și O(A·B) memorie.
Pentru N ≤ 199 aceasta este de ordinul a câtorva zeci de mii de operaţii.

DP pentru submulţimile fără repetiție – O(M·A·B) timp și O(A·B) memorie (se poate reduce la două straturi pentru a economisi spaţiu).
În cel mai nefavorabil caz M = N = 199, A și B nu depășesc 199, deci numărul de operaţii este sub 8·10⁶, confortabil în limita de timp.

Calculul final – O(K) timp.

În total, algoritmul rulează în timp O(N³) cu un factor mic și consumă O(N²) memorie, mult sub limitele impuse de problemă.

9. Concluzie
Prin descompunerea procesului de extragere în (a) distribuţia numărului de şosete mici/mari și (b) numărarea combinărilor valide de culori pentru fiecare repartizare, am transformat problema probabilistică într‑o serie de calcule combinatorii exacte. Lemma 3 garantează că DP‑ul pentru culori numără corect submulţimile fără repetiție, iar normalizarea cu combinaţiile totale transformă acest număr în probabilitatea dorită. Îmbinarea cu distribuţia q duce la formula finală, care este exactă și poate fi evaluată cu precizia cerută (eroare absolută ≤ 10⁻⁵). Astfel soluţia propusă este corectă, eficientă și ușor de implementat.